% !PRINTED 269
\weberpage{269}
%<<CONT
substitution of the third kind from $\Lambda_p$, then we can determine
the substitutions $U$, $V$ from $L_p$ or from $\Lambda_p$ in such a way that
\[
A = U\,S_2^r\,U^{-1}, \qquad A = V\,S_3^r\,V^{-1}.
\]
If, then, corresponding to the value $r=1$, we set
\[
A_1 = U\,S_2\,U^{-1}, \qquad A_1 = V\,S_3\,V^{-1},
\]
it follows that
\[
A = A_1^r, \qquad A = A_1^r,
\]
and $A_1$ occurs in $L_p$, $A_1$ in $\Lambda_p$.

If, further, $B$ is a substitution of the first kind from $L_p$, then,
setting
\[
T_1 = \begin{pmatrix} 1, & 1 \\ 0, & 1 \end{pmatrix}
\]
one can again determine $U$ from $L_p$ in such a way that
\[
B = U\,T_1^r\,U^{-1} = B_1^r
\]
holds. From this results the theorem:

\gesperrt{In the group $L_p$ (or $\Lambda_p$) the substitutions
of the first, second and third kinds can, with the inclusion
of the identity, be arranged in cycles of $p$, $\dfrac{p-1}{2}$, $\dfrac{p+1}{2}$
terms.}

Two such cycles can have no term in common besides the
unit element.

This is at once evident for the substitutions of the first
kind, whose degree is equal to $p$; for with these the whole
cycle can be derived by raising to powers from any element
different from unity.

But if a common element occurs in two cycles of the
second or third kind, then the group can be so transformed
that the two cycles take the form:
\weq{12}{\begin{aligned} &1, \quad U\,S\,U^{-1}, \quad U\,S^2\,U^{-1}, \ \ldots \\ &1, \quad S, \quad S^2, \ \ldots, \end{aligned}}
in which $S = \begin{pmatrix} \sigma, & 0 \\ 0, & \sigma^{-1} \end{pmatrix}$ has the normal form.

If now, for any two exponents
$r,t$ different from zero,
\[
S^r = U\,S^t\,U^{-1},
\]
%CONT>>
